\documentclass[11pt]{article}
\usepackage[letterpaper,margin=1in]{geometry}
\usepackage{amsmath,amssymb,booktabs,hyperref}
\hypersetup{hidelinks,pdftitle={A two-shift bound for the least quadratic nonresidue}}
\newcommand{\paperAuthor}{Lennart Lopin\\ {\normalsize Euler's Identity LLC}\\ {\normalsize ORCID: \href{https://orcid.org/0009-0001-2751-2060}{0009-0001-2751-2060}}}
% Enter the public release DOI URL or a permanent repository/commit URL here.
% Filling this one field enables the in-text software citation and bibliography item.
\newcommand*{\replicationURL}{https://doi.org/10.5281/zenodo.24408857}
\newcommand*{\ifReplicationURL}[1]{\ifx\replicationURL\empty\else#1\fi}
\title{A two-shift bound for the least quadratic nonresidue}
\author{\paperAuthor}
\date{November 20, 2026}
\newcommand{\CC}{\mathcal C(1)}
\newcommand{\abs}[1]{\left|#1\right|}
\newcommand{\RR}{\mathbb R}
\newcommand{\ii}{\mathrm i}
\begin{document}
\maketitle
\begin{abstract}
We give an explicit two-shift Fourier test function certifying
$\mathcal C(1)^{-2}<0.76079030$ for the extremal constant of Carneiro,
Milinovich, Quesada-Herrera and Ramos. Their published upper bound is
$0.7615$. The witness has 24 rational coefficients, and three separately written interval
computations certify its performance. Under the generalized Riemann
hypothesis for Dirichlet $L$-functions, their Theorems 1 and 2 give the
same improved upper bound for the limiting ratios of the least quadratic
nonresidue and the least prime quadratic residue to $(\log p)^2$.
The Fourier inequality itself is unconditional.
\end{abstract}

\section{The bound and its two consequences}
For an odd prime $p$, let $n_p$ be the least positive quadratic nonresidue
modulo $p$ and $r_p$ the least prime that is a nonzero quadratic residue
modulo $p$. Write $\widehat F(t)=\int_{\RR} F(x)e^{-2\pi\ii xt}\,dx$. For
nonzero $F\in L^1(\RR)$ with real-valued transform, define $D=\|F\|_1$ and
\[
 N=2\pi\left(\int_{-\infty}^0 \widehat F(t)e^{\pi t}\,dt
                 -\int_0^\infty |\widehat F(t)|e^{\pi t}\,dt\right).
\]
The supremum of $N/D$ is $\CC$ in \cite[(1.1)]{cmqr}. Each admissible
witness with $N>0$ therefore gives $\CC^{-2}\le(D/N)^2$.

\medskip\noindent
\textbf{Proposition.} The function defined in Section 2 satisfies
$(D/N)^2<0.76079030$. Consequently, assuming Dirichlet GRH,
\[
 \max\!\left\{\limsup_{p\to\infty}\frac{n_p}{(\log p)^2},
               \limsup_{p\to\infty}\frac{r_p}{(\log p)^2}\right\}
 \le \CC^{-2}<0.76079030 ,
\]
where $p$ runs over primes. The first inequality is precisely the joint
consequence of Theorems 1 and 2 of \cite{cmqr}. The improvement is in the
coefficient of $(\log p)^2$. The authors also give the method-specific
lower bound $0.7596<\CC^{-2}$.

Alp\"oge's announcement of a Lean formalization covering Dirichlet
GRH~\cite{alpoge} motivated the search for this witness. Sutherland reports
independent checks of the formal statements~\cite{sutherland}. Taking the
Dirichlet GRH result as established discharges the hypothesis of both
arithmetic corollaries above.

\newpage
\section{The witness}
Take $a=0.246$, $c_1=0.85$, $c_2=0$, and use the integers in Table 1 to
define $b_{jn}=B_{jn}/10^7$. Put
\[
 P_j(u)=\sum_{n=1}^{12} b_{jn}\frac{u^{2n-1}}{(2n-1)!},\qquad
 u_j=\frac{c_j-\pi t}{a},\qquad
 \widehat F(t)=\sum_{j:\,u_j\ge0}e^{-u_j}P_j(u_j).
\]
Fourier inversion and the gamma integral give
\[
 F(x)=\frac a\pi\sum_{j,n} b_{jn} e^{2\ii c_jx}(1+2\ii ax)^{-2n}.
\]
Thus $F\in L^1$, and its Fourier transform is real and continuous. The
two translations allow the shape on negative frequencies to change while
only the first component reaches positive frequencies.

\begin{table}[h]
\centering
\begin{tabular}{r rr@{\qquad}r rr}
\toprule
$n$ & $B_{1n}$ & $B_{2n}$ & $n$ & $B_{1n}$ & $B_{2n}$ \\
\midrule
1 & 1295 & 0 & 7 & -46679113 & -118505260 \\
2 & 79929 & -9074 & 8 & 20002165 & 154829818 \\
3 & -6609595 & 5043626 & 9 & -2996805 & -138162493 \\
4 & 117992958 & -14383101 & 10 & -22963877 & 320107061 \\
5 & -234402020 & -118792841 & 11 & 279853645 & 234086697 \\
6 & 332397224 & 61384470 & 12 & -809901785 & 62986550 \\
\bottomrule
\end{tabular}
\caption{Exact integer coefficients in ascending odd powers.}
\end{table}

\paragraph{Numerator.} Set $w=\pi t$ and $\lambda=1+a$. Integrating the
signed transform over all $w$ and subtracting twice its positive part on
$w>0$ yields
\[
 N=2a\sum_{j,n}\frac{e^{c_j}b_{jn}}{(1+a)^{2n}}
       -4ae^{c_1}\int_0^{c_1/a}e^{-\lambda u}\max\{P_1(u),0\}\,du .
\]
Here $\int_0^\infty e^{-\lambda u}u^{2n-1}/(2n-1)!\,du=\lambda^{-2n}$.

\paragraph{Norm.} After removing a common phase, let $z=2ax$,
$\omega_j=(c_j-c_1)/a$, and put
\[
 H(z)=\sum_{j,n} b_{jn} e^{\ii\omega_j z}(1+\ii z)^{-2n},\qquad
 D=\frac1\pi\int_0^\infty |H(z)|\,dz .
\]
The factor $1/\pi$ uses $H(-z)=\overline{H(z)}$ for real $z$. For $R>0$,
the triangle inequality and $|1+\ii z|\ge z$ for $z>0$ give
\begin{equation}
 \int_R^\infty |H(z)|\,dz
 \le\sum_{j,n}\frac{|b_{jn}|}{(2n-1)R^{2n-1}} . \label{eq:tail}
\end{equation}

\newpage
\section{Computational certificate}
All three verifiers use the exact decimals in \texttt{witness.json} and
the transform identities in Section 2, and bound the tail by
\eqref{eq:tail}. They use different finite integration routines.
The search that produced the coefficients is unnecessary for verification.

\paragraph{Complex ball quadrature.}
The Python-FLINT verifier uses Arb at 128-bit precision on finite panels.
For $\epsilon=10^{-12}$, it bounds $2\max(y,0)$ by
$y+\sqrt{y^2+\epsilon^2}$ in the numerator correction. For the norm it
integrates $\sqrt{H(z)H(-z)+\epsilon^2}$, an upper bound on the real
axis. Both square roots use Arb's analyticity check on complex boxes.
The resulting outward-rounded inequalities are
\[
 N>0.9177970191817,\qquad D<0.8005327903879 .
\]
These imply $(D/N)^2<0.76079030$ and prove the proposition.

\paragraph{Trapezoid check.}
The mpmath verifier encloses endpoint values at 35 decimal digits. On a
panel of length $h$, linear interpolation of a real or complex $g$ has
pointwise error at position $s$ at most $s(h-s)\sup|g''|/2$.
Convexity and the Lipschitz constant 1 of modulus and positive part
therefore give a trapezoid upper bound with additive error
$h^3\sup|g''|/12$, even on panels containing a zero. The code bounds
the derivatives term by term. This check gives $(D/N)^2<0.760799446$.

\paragraph{Simpson check.}
A third verifier shares no code with the first two. It isolates the two
sign changes of $P_1$ on $(0,c_1/a)$ by exact rational root isolation and
integrates the positive part in closed form. For the norm it applies
adaptive Simpson panels on $[0,96]$, bounding the fourth derivative of
$|H|$ on each panel by Taylor-ball arithmetic at 150 bits, and uses
\eqref{eq:tail} with $R=96$. This check gives $(D/N)^2<0.760790289$.

\paragraph{Lean verification.}
The accompanying Lean 4 file proves the infinite-tail estimate
\eqref{eq:tail} for finite sums of the indicated type, checks that its right side at $R=64$
is below $0.000002035$ for Table 1, and proves the final
ratio inequality assuming the stated numerical enclosures for $N$ and $D$.
It compiles with Lean and Mathlib v4.24.0 with warnings treated as errors,
and these declarations depend only on the axioms \texttt{propext},
\texttt{Classical.choice} and \texttt{Quot.sound}.
The Python interval verifiers supply the enclosures. The Fourier reduction
is derived in Section 2, and the number-theoretic implications use
\cite{cmqr}.

\paragraph{Reproduction and provenance.}
Each of the following commands, run in the archive directory, checks the witness:
\begin{itemize}
\item \texttt{python3 check\_arb.py witness.json} (python-flint 0.9.0);
\item \texttt{python3 certify.py witness.json} (mpmath 1.3.0);
\item \texttt{python3 verify\_two\_shift\_simpson.py witness.json}\newline
      (python-flint 0.9.0 and SymPy 1.14.0).
\end{itemize}
The companion archive includes the input, the three verifiers, optimization
code and Lean source\ifReplicationURL{~\cite{replication}}. OpenAI Codex assisted with the numerical search,
verification code and preparation of this note. Anthropic's Claude
independently reran all three verifiers and the Lean check before deposit.

\begin{thebibliography}{1}
\bibitem{cmqr} E. Carneiro, M. B. Milinovich, E. Quesada-Herrera and
A. P. Ramos, \emph{Fourier optimization, the least quadratic non-residue,
and the least prime in an arithmetic progression}, Mathematics of
Computation, \href{https://doi.org/10.1090/mcom/4154}{doi:10.1090/mcom/4154}.
\href{https://arxiv.org/abs/2404.08380}{arXiv:2404.08380}.
\bibitem{alpoge} L.~Alp\"oge, [Announcement of a Lean formalization of the
Riemann hypothesis], November 20, 2026.
\url{https://alpo.ge/remark.pdf}.

\bibitem{sutherland} A.~V.~Sutherland, \emph{Verification files for the
Lean proof of the Riemann Hypothesis}, MIT, November 20, 2026.
\url{https://math.mit.edu/~drew/grh/}.

\ifReplicationURL{\bibitem{replication} \emph{Replication materials for A
two-shift bound for the least quadratic nonresidue}. Exact coefficients,
verification code, and Lean source.
\expandafter\url\expandafter{\replicationURL}. Source repository:
\url{https://github.com/novalis78/least-nonresidue-two-shift}.}
\end{thebibliography}
\end{document}
